% !TeX program = pdflatex
% Architektura zálohování a ochrany dat aplikace Tappka
% Sazba:  latexmk -pdf zalohovani-tappka.tex
%   nebo: pdflatex zalohovani-tappka.tex  (dvakrát kvůli obsahu a odkazům)

\documentclass[11pt,a4paper]{article}

\usepackage[T1]{fontenc}
\usepackage[utf8]{inputenc}
\usepackage[czech]{babel}
\usepackage{lmodern}
\usepackage{amssymb}
\usepackage[margin=2.4cm,headheight=14pt]{geometry}

\usepackage{booktabs}
\usepackage{tabularx}
\usepackage{longtable}
\usepackage{array}
\usepackage{multirow}
\usepackage{pifont}
\usepackage{enumitem}
\usepackage{xcolor}
\usepackage{tcolorbox}
\usepackage{listings}
\usepackage{fancyhdr}
\usepackage{microtype}
\usepackage[unicode,hidelinks]{hyperref}

% --- Barvy ---------------------------------------------------------------
\definecolor{accent}{HTML}{1F4E79}
\definecolor{warnbg}{HTML}{FDF3F3}
\definecolor{warnln}{HTML}{C0392B}
\definecolor{infobg}{HTML}{F1F6FB}
\definecolor{codebg}{HTML}{F7F7F7}
\definecolor{muted}{HTML}{5A6672}

% --- Značky --------------------------------------------------------------
\newcommand{\ano}{\textcolor{accent}{\ding{51}}}
\newcommand{\nepokr}{\textcolor{warnln}{\ding{55}}}
\newcommand{\castecne}{\textcolor{muted}{\ding{108}}}

% --- Hlavička a patička --------------------------------------------------
\pagestyle{fancy}
\fancyhf{}
\fancyhead[L]{\small\textcolor{muted}{Architektura zálohování aplikace Tappka}}
\fancyhead[R]{\small\textcolor{muted}{verze 1.0}}
\fancyfoot[C]{\small\thepage}
\renewcommand{\headrulewidth}{0.4pt}

% --- Nadpisy -------------------------------------------------------------
\usepackage{titlesec}
\titleformat{\section}{\Large\bfseries\color{accent}}{\thesection}{0.7em}{}
\titleformat{\subsection}{\large\bfseries}{\thesubsection}{0.6em}{}

% --- Boxy ----------------------------------------------------------------
\newtcolorbox{upozorneni}[1][]{
  colback=warnbg, colframe=warnln, boxrule=0.9pt, arc=2pt,
  left=8pt, right=8pt, top=6pt, bottom=6pt,
  fonttitle=\bfseries, title={#1}}

\newtcolorbox{informace}[1][]{
  colback=infobg, colframe=accent, boxrule=0.9pt, arc=2pt,
  left=8pt, right=8pt, top=6pt, bottom=6pt,
  fonttitle=\bfseries, title={#1}}

% --- Výpisy kódu ---------------------------------------------------------
\lstdefinestyle{konfig}{
  basicstyle=\ttfamily\footnotesize,
  backgroundcolor=\color{codebg},
  frame=single, rulecolor=\color{muted!40},
  breaklines=true, breakatwhitespace=true,
  showstringspaces=false,
  keywordstyle=\color{accent}\bfseries,
  commentstyle=\color{muted}\itshape,
  numbers=none, xleftmargin=6pt, xrightmargin=2pt,
  aboveskip=10pt, belowskip=10pt,
  columns=fullflexible
}
\lstset{style=konfig}

\setlist{itemsep=2pt, topsep=4pt}
\renewcommand{\arraystretch}{1.25}
\setlength{\emergencystretch}{3em}

% =========================================================================
\begin{document}

\begin{titlepage}
\centering
\vspace*{2.5cm}

{\Huge\bfseries\color{accent} Architektura zálohování\\[0.35em] a ochrany dat aplikace Tappka\par}

\vspace{1.2cm}
{\Large Návrh třívrstvého modelu ochrany dat\par}
\vspace{0.4cm}
{\large\color{muted} Podklad pro jednání o alokaci serverových kapacit\par}

\vspace{3cm}

\begin{tabular}{@{}ll@{}}
\textbf{Verze dokumentu} & 1.0 \\
\textbf{Datum} & 31.\,8.\,2026 \\
\textbf{Zpracoval:a} & Vývojový tým Tappka \\
\textbf{Určeno} & Vedení školy, IT oddělení, oddělení PEVKA \\
\textbf{Stav} & Návrh k projednání \\
\end{tabular}

\vfill
{\small\color{muted} Všechna technická tvrzení o platformě Supabase v tomto dokumentu jsou doložena\\ odkazy na oficiální dokumentaci v kapitole \ref{sec:zdroje} (Zdroje).\par}
\end{titlepage}

\tableofcontents
\newpage

% =========================================================================
\section{Shrnutí pro vedení}

Tappka je webová aplikace, ve které vzniká a je uložena podstatná část studijních a reflektivních
záznamů studujících. Ztráta těchto dat by nebyla obnovitelná z jiného zdroje --- neexistuje
papírová ani jiná paralelní evidence.

Aplikace dnes běží na platformě \textbf{Supabase} v tarifu \textbf{Pro}. Tento tarif zahrnuje
automatické denní zálohy databáze s~retencí sedmi dní. To je solidní základ, ale jako celková
ochrana dat \textbf{nestačí}, a to ze tří konkrétních důvodů:

\begin{enumerate}[leftmargin=1.6em]
  \item \textbf{Zálohy Supabase nezahrnují nahrané soubory.} Zálohuje se pouze databáze.
        Fotografie, profilové obrázky a nahrané dokumenty leží v odděleném úložišti, které
        automatické zálohy \emph{nepokrývají vůbec}. Smazaný soubor je nenávratně ztracen
        i~v~tarifu Pro (viz kap.~\ref{sec:storage-nezalohy}).
  \item \textbf{Sedm dní je krátká lhůta.} Problém zjištěný po prázdninách, po zkouškovém období
        nebo po delší nepřítomnosti vyučující:ho už není z čeho obnovit.
  \item \textbf{Všechny kopie leží u jednoho dodavatele.} Zálohy Supabase jsou uloženy v~téže
        infrastruktuře jako produkční data. Nechrání proti výpadku, ukončení služby ani ztrátě
        přístupu k~účtu.
\end{enumerate}

\begin{informace}[Co navrhujeme]
Doplnit stávající ochranu o~\textbf{dvě nezávislé vrstvy}, které společně tvoří třívrstvý model:

\smallskip
\begin{tabularx}{\textwidth}{@{}p{2.6cm}X@{}}
\textbf{Vrstva 1} & Platformní ochrana Supabase Pro --- denní zálohy databáze, retence 7~dní.
                    Již funguje, beze změny nákladů. \\
\textbf{Vrstva 2} & \emph{Nezávislá záloha na serveru školy.} Kontejnerizovaná služba, která
                    jednou týdně stáhne kompletní kopii databáze \emph{i~všech nahraných souborů}
                    a~uloží ji šifrovaně na školní úložiště. \textbf{Toto je hlavní předmět
                    tohoto dokumentu.} \\
\textbf{Vrstva 3} & Lidsky čitelné exporty --- vize do budoucna. Data v~podobě, kterou lze
                    otevřít a~přečíst i~bez Tappky a~bez technických znalostí. \\
\end{tabularx}
\end{informace}

\subsection*{O co konkrétně žádáme školu}

\begin{center}
\begin{tabularx}{\textwidth}{@{}p{4.6cm}X@{}}
\toprule
\textbf{Položka} & \textbf{Specifikace} \\
\midrule
Výpočetní kapacita & Virtuální server, Linux (Debian 12 nebo Ubuntu 24.04 LTS), 2~vCPU, 4~GB~RAM \\
Diskový prostor & \textbf{1~TB} dedikovaného prostoru (odůvodnění kalkulací v~kap.~\ref{sec:kapacita}) \\
Software & Docker Engine s~pluginem Docker Compose \\
Síť & Odchozí HTTPS (TCP 443) a~TCP 5432 na infrastrukturu Supabase \\
Provozní součinnost & Předání zálohy na vyžádání v~dohodnutých lhůtách (kap.~\ref{sec:vyzadani}) \\
\bottomrule
\end{tabularx}
\end{center}

\medskip
Řešení je navrženo tak, aby po jednorázovém uvedení do provozu \textbf{nevyžadovalo průběžnou
správu} --- běží automaticky a~upozorní pouze v~případě selhání.

\newpage

% =========================================================================
\section{Kontext a rozsah}

\subsection{Jaká data Tappka spravuje}

Z pohledu zálohování je podstatné, že data Tappky nejsou homogenní. Dělí se do kategorií, které se
liší mírou nenahraditelnosti:

\begin{center}
\begin{tabularx}{\textwidth}{@{}p{4.0cm}X>{\raggedright\arraybackslash}p{3.1cm}@{}}
\toprule
\textbf{Kategorie} & \textbf{Příklad obsahu} & \textbf{Nahraditelnost} \\
\midrule
Autorský obsah studujících & Eseje, reflexe, sebehodnocení, roční týmové reflexe & Nenahraditelné \\
Hodnocení a~zpětná vazba & Komentáře kouči:ek, výsledky, bodové zisky & Nenahraditelné \\
Procesní záznamy & Rezervace, docházka na aktivity, výpůjčky knih & Obtížně rekonstruovatelné \\
Nahrané soubory & Profilové fotografie, dokumenty, obrázky v~esejích & Nenahraditelné \\
Konfigurace a~číselníky & Týmy, role, tagy, kategorie, nastavení & Rekonstruovatelné \\
\bottomrule
\end{tabularx}
\end{center}

Většina objemu dat připadá na poslední dvě kategorie, ale \emph{většina hodnoty} na první tři.
Zálohovací strategie proto musí pokrýt obojí --- nestačí zálohovat jen to, co zabírá místo.

\subsection{Dva nezávislé systémy uložení}
\label{sec:dva-systemy}

Klíčový fakt, na kterém stojí celý zbytek dokumentu: \textbf{Tappka ukládá data do dvou technicky
oddělených systémů.}

\begin{center}
\begin{tabularx}{\textwidth}{@{}p{3.5cm}p{4.2cm}X@{}}
\toprule
\textbf{Systém} & \textbf{Co obsahuje} & \textbf{Technologie} \\
\midrule
\textbf{Databáze} & Texty, hodnocení, vztahy, metadata & PostgreSQL \\
\textbf{Úložiště souborů} & Binární soubory --- fotografie, dokumenty, obrázky & Objektové úložiště kompatibilní s~S3, spravované samostatnou službou \\
\bottomrule
\end{tabularx}
\end{center}

Aplikace používá tři oddělené kontejnery (\emph{buckety}) v~úložišti souborů:

\begin{itemize}
  \item \texttt{avatars} --- profilové obrázky osob a~týmů (veřejně čitelné)
  \item \texttt{images} --- obálky knih a~obrázky vložené do esejí (veřejně čitelné)
  \item \texttt{documents} --- nahrané dokumenty (neveřejné, přístupné jen přes dočasné
        podepsané odkazy)
\end{itemize}

V databázi jsou o~těchto souborech uloženy pouze \emph{odkazy a~metadata} --- název, velikost,
kdo soubor nahrál. Samotný obsah souboru v~databázi není. Tento rozdíl je příčinou nejzávažnější
mezery v~současné ochraně, popsané v~kapitole~\ref{sec:storage-nezalohy}.

\subsection{Kde jsou data fyzicky uložena}
\label{sec:lokalita}

\begin{center}
\begin{tabularx}{\textwidth}{@{}p{4.4cm}X@{}}
\toprule
\textbf{Parametr} & \textbf{Hodnota} \\
\midrule
Region & Central Europe (Zurich) \\
Označení regionu & \texttt{eu-central-2} \\
Stát & Švýcarská konfederace \\
Poskytovatel infrastruktury & Amazon Web Services (Supabase provozuje na AWS) \\
Výpočetní instance & \texttt{t4g.micro} (tarif Micro: 2~jádra ARM, 1~GB~RAM, limit 60 souběžných připojení) \\
Zálohy & Uloženy v~téže infrastruktuře a~témže regionu jako produkční databáze \\
\bottomrule
\end{tabularx}
\end{center}

\subsubsection*{Právní kontext přenosu dat mimo EU}

Švýcarsko není členem Evropské unie ani EHP. Přenos osobních údajů do Švýcarska je nicméně
z~pohledu GDPR postaven na roveň přenosu uvnitř EU: Evropská komise dne \textbf{15.~1.~2024}
potvrdila, že Švýcarsko nadále zajišťuje odpovídající úroveň ochrany osobních údajů. Osobní údaje
lze tedy z~EU do Švýcarska předávat \textbf{bez nutnosti dalších záruk} (standardních smluvních
doložek apod.) --- viz zdroj~\cite{adekvatnost}.

Volba curyšského regionu je z~hlediska ochrany dat příznivá: data neopouštějí evropský prostor
a~zůstávají geograficky blízko uživatel:ek.

\begin{informace}[Dopad na návrh vrstvy 2]
Skutečnost, že produkční data i~jejich zálohy leží v~jednom regionu u~jednoho poskytovatele, je
hlavní argument pro vrstvu 2. Záloha uložená vedle originálu nechrání proti výpadku
poskytovatele, ztrátě přístupu k~účtu ani ukončení služby. Umístění kopie \textbf{na serveru
školy v~České republice} vytváří skutečnou geografickou i~organizační nezávislost.
\end{informace}

\newpage

% =========================================================================
\section{Analýza rizik}
\label{sec:rizika}

Zálohování není samoúčelné --- chrání proti konkrétním scénářům. Následující tabulka je zároveň
odůvodněním, proč jedna vrstva nestačí.

\begin{center}
\begin{tabularx}{\textwidth}{@{}Xccc@{}}
\toprule
\textbf{Scénář} & \textbf{Vrstva 1} & \textbf{Vrstva 2} & \textbf{Vrstva 3} \\
 & {\footnotesize Supabase} & {\footnotesize školní server} & {\footnotesize exporty} \\
\midrule
Omylem smazaný záznam, zjištěno do 7 dní & \ano & \ano & \castecne \\
Omylem smazaný nebo přepsaný \textbf{soubor} & \nepokr & \ano & \ano \\
Chyba v~migraci, poškození struktury databáze & \ano & \ano & \nepokr \\
Problém zjištěný \textbf{po více než 7 dnech} & \nepokr & \ano & \ano \\
Kompromitace účtu, hromadné smazání dat & \castecne & \ano & \ano \\
Delší výpadek platformy Supabase & \nepokr & \ano & \ano \\
Ukončení služby Supabase nebo sporu s~dodavatelem & \nepokr & \ano & \ano \\
Ukončení provozu aplikace Tappka & \nepokr & \castecne & \ano \\
Potřeba dat pro osobu bez technických znalostí & \nepokr & \nepokr & \ano \\
Selhání či ztráta školního serveru & \ano & \nepokr & \castecne \\
\bottomrule
\end{tabularx}
\end{center}

\medskip
\noindent\small
\ano{}~pokrývá \quad \castecne{}~pokrývá částečně \quad \nepokr{}~nepokrývá
\normalsize

\medskip
Poslední řádek je záměrný: vrstvy se navzájem jistí. Neexistuje jediná událost, která by při
zavedení všech tří vrstev vedla ke ztrátě dat.

\subsection{Odborný rámec: pravidlo 3--2--1}

Navržená architektura odpovídá pravidlu \textbf{3--2--1}, které jako základní standard ochrany dat
doporučují americká agentura CISA i~institut NIST~\cite{pravidlo321}:

\begin{center}
\begin{tabularx}{\textwidth}{@{}p{3.8cm}X@{}}
\toprule
\textbf{Požadavek} & \textbf{Naplnění v~návrhu} \\
\midrule
\textbf{3} kopie dat & Produkční data (Supabase) $+$ zálohy Supabase $+$ kopie na školním serveru \\
\textbf{2} různá média & Spravovaná cloudová infrastruktura AWS $+$ diskové pole školy \\
\textbf{1} kopie mimo lokalitu & Školní server v~ČR, jiný poskytovatel, jiná organizace, jiný stát \\
\bottomrule
\end{tabularx}
\end{center}

Moderní rozšíření pravidla, \textbf{3--2--1--1--0}, přidává jednu neměnitelnou (immutable) kopii
a~\emph{nulu chyb}, tedy pravidelně \textbf{ověřovanou} obnovu. Oba tyto prvky jsou v~návrhu
obsaženy --- viz kapitoly \ref{sec:zabezpeceni} a~\ref{sec:overovani}.

\newpage

% =========================================================================
\section{Vrstva 1 --- ochrana na platformě Supabase Pro}
\label{sec:vrstva1}

Tato vrstva je již v~provozu a~nevyžaduje žádnou akci ze strany školy. Je zde popsána proto, aby
bylo zřejmé, \emph{co přesně pokrývá a~co nikoli}.

\subsection{Automatické denní zálohy}

\begin{center}
\begin{tabularx}{\textwidth}{@{}p{4.4cm}X@{}}
\toprule
\textbf{Parametr} & \textbf{Hodnota} \\
\midrule
Frekvence & Jednou denně, automaticky \\
Retence & \textbf{7 dní} (posledních sedm denních záloh) \\
Typ zálohy & Fyzická záloha (snímek diskového svazku PostgreSQL) \\
Rozsah & Databáze PostgreSQL \\
Obnova & Samoobslužně z~administračního rozhraní (\emph{Database~$\rightarrow$~Backups}) nebo přes API \\
Nejhorší možná ztráta dat (RPO) & Až 24 hodin \\
Doba obnovy (RTO) & Řádově desítky minut až hodiny podle velikosti databáze \\
Dostupnost aplikace při obnově & \textbf{Nedostupná} --- projekt je po dobu obnovy odstaven \\
\bottomrule
\end{tabularx}
\end{center}

Zdroj: oficiální dokumentace Supabase~\cite{zalohysupabase} a~ceník~\cite{cenik}.

\subsection{Co zálohy Supabase nepokrývají}
\label{sec:storage-nezalohy}

\begin{upozorneni}[Zásadní omezení: nahrané soubory nejsou zálohovány]
Dokumentace Supabase uvádí doslova, že zálohy databáze \emph{nezahrnují objekty uložené přes
Storage API} --- databáze o~nich obsahuje pouze metadata~\cite{zalohysupabase}.

\smallskip
Prakticky to znamená: pokud dojde ke smazání nebo poškození souboru v~kontejnerech
\texttt{avatars}, \texttt{images} nebo \texttt{documents}, \textbf{neexistuje žádná záloha, ze
které by šel obnovit.} Denní zálohy obnoví pouze databázový záznam odkazující na soubor, který už
neexistuje --- výsledkem je nefunkční odkaz.

\smallskip
Toto omezení není vadou konkrétního tarifu. Vyplývá z~toho, že zálohy jsou fyzické snímky
diskového svazku databáze, zatímco soubory leží v~samostatném objektovém úložišti mimo tento
svazek. Omezení tedy platí ve všech tarifech.
\end{upozorneni}

Kromě souborů se zálohy netýkají:
\begin{itemize}
  \item hesel vlastních databázových rolí (nejsou v~denních zálohách obsažena),
  \item trvale smazaných objektů (\emph{deleted objects cannot be recovered}),
  \item replikačních slotů a~odběrů, které je nutné po obnově ručně obnovit.
\end{itemize}

\subsection{Bezpečnostní opatření platformy}

Tarif Pro poskytuje následující bezpečnostní prvky, které Tappka využívá nebo může využít:

\begin{center}
\begin{tabularx}{\textwidth}{@{}p{4.6cm}Xc@{}}
\toprule
\textbf{Opatření} & \textbf{Popis} & \textbf{Stav} \\
\midrule
Zabezpečení na úrovni řádků (RLS) & Přístupová pravidla vynucovaná přímo databází, nikoli aplikací. Uživatel:ka nemůže obejít pravidla ani přímým dotazem na API. & \ano zavedeno \\
Vynucení TLS (\emph{SSL Enforcement}) & Odmítnutí nešifrovaných databázových spojení & doporučeno \\
Síťová omezení & Omezení přístupu k~databázi na povolené rozsahy IP adres & doporučeno \\
Vícefaktorové ověření & Ochrana účtů se správcovským přístupem k~projektu & \ano zavedeno \\
Omezení četnosti požadavků & Ochrana přihlašování proti automatizovaným útokům & \ano výchozí \\
Retence provozních záznamů & \textbf{7 dní} historie logů & součást tarifu \\
Šifrování přenosu & TLS mezi aplikací, databází i~úložištěm souborů & \ano výchozí \\
\bottomrule
\end{tabularx}
\end{center}

Zdroj: produkční kontrolní seznam Supabase~\cite{prodchecklist} a~ceník~\cite{cenik}.

\subsection{Co tarif Pro neobsahuje}

Pro úplnost a~korektnost podkladu je nutné uvést i~to, co v~tarifu Pro chybí:

\begin{itemize}
  \item \textbf{Přístup k~auditní zprávě SOC 2.} Společnost Supabase je jako organizace držitelem
        certifikace SOC~2 Type~2 a~je každoročně auditována~\cite{soc2}. Samotnou auditní zprávu
        si však mohou stáhnout pouze zákazníci v~tarifech Team a~Enterprise. V~tarifu Pro tedy
        \emph{platforma je certifikovaná, ale doklad o~tom není k~dispozici.}
  \item \textbf{Certifikace ISO~27001 a~smlouva o~zpracování dle HIPAA} --- nabízeny až ve vyšších
        tarifech.
  \item \textbf{Smluvní garance dostupnosti (SLA).} Tarif Pro zahrnuje e-mailovou podporu, nikoli
        závazek dostupnosti služby~\cite{cenik}.
  \item \textbf{Jednotné přihlášení (SSO) pro správu organizace.}
\end{itemize}

\begin{informace}[Praktický dopad]
Absence SLA znamená, že v~případě výpadku Supabase \textbf{neexistuje smluvně vymahatelná lhůta
obnovení služby}. Škola je odkázána na dobrou vůli dodavatele. Vrstva 2 tuto závislost odstraňuje:
i~při delším výpadku zůstávají data v~držení školy a~lze je použít.
\end{informace}

\subsection{Model sdílené odpovědnosti}

Supabase publikuje explicitní model sdílené odpovědnosti~\cite{sdilenaodpovednost}. Zjednodušeně:

\begin{center}
\begin{tabularx}{\textwidth}{@{}p{5.2cm}X@{}}
\toprule
\textbf{Odpovědnost Supabase} & \textbf{Odpovědnost provozovatele aplikace} \\
\midrule
Provoz a~zálohování databáze, údržba operačního systému, správa a~monitoring infrastruktury &
Návrh datového modelu, nastavení přístupových práv a~RLS, správa tajemství, \textbf{vlastní
zálohovací strategie nad rámec platformy} \\
\bottomrule
\end{tabularx}
\end{center}

Dokumentace to shrnuje větou, kterou je vhodné citovat: \emph{,,Supabase klade kolem vaší databáze
jen velmi málo zábran. To vám dává velkou kontrolu, ale zároveň to znamená, že věci můžete
rozbít.``}~\cite{sdilenaodpovednost} Odpovědnost za data tedy zůstává na straně provozovatele ---
a~tím je v~konečném důsledku škola.

\newpage

% =========================================================================
\section{Vrstva 2 --- nezávislá záloha na serveru školy}
\label{sec:vrstva2}

Toto je jádro návrhu a~hlavní důvod vzniku tohoto dokumentu.

\subsection{Cíl}

Vytvořit a~automaticky udržovat \textbf{úplnou, samostatně použitelnou kopii} všech dat Tappky
--- databáze \emph{i} nahraných souborů --- na infrastruktuře, kterou vlastní a~kontroluje škola,
nezávisle na dodavateli.

\subsection{Architektura}

Řešení je navrženo jako jediná kontejnerizovaná služba spouštěná plánovačem \texttt{cron}.
Kontejnerizace (Docker) je zvolena záměrně, ze tří důvodů:

\begin{enumerate}[leftmargin=1.6em]
  \item \textbf{Nulové nároky na prostředí serveru.} Kontejner nese vlastní verze všech nástrojů.
        Škola neinstaluje PostgreSQL, Python ani nic dalšího --- pouze Docker.
  \item \textbf{Shoda verzí.} Nástroj \texttt{pg\_dump} musí odpovídat verzi serveru PostgreSQL.
        Kontejner tuto shodu zaručuje a~izoluje ji od aktualizací školního systému.
  \item \textbf{Reprodukovatelnost a~přenositelnost.} Stejný kontejner lze spustit na jiném
        stroji a~ověřit obnovu, aniž by se cokoli měnilo v~konfiguraci.
\end{enumerate}

\subsubsection*{Průběh jednoho zálohovacího cyklu}

\begin{center}
\begin{tabularx}{\textwidth}{@{}p{0.6cm}p{4.0cm}X@{}}
\toprule
\textbf{\#} & \textbf{Krok} & \textbf{Popis} \\
\midrule
1 & Spuštění & \texttt{cron} v~kontejneru spustí zálohovací skript (neděle 02:00) \\
2 & Výpis databáze & Tři samostatné soubory: role, struktura, data \\
3 & Kontrola & Ověření, že žádný z~výpisů není prázdný; při chybě cyklus končí a~upozorní \\
4 & Synchronizace souborů & Přírůstkové stažení všech tří kontejnerů úložiště \\
5 & Zabalení a~šifrování & Komprese a~zašifrování archivu veřejným klíčem \\
6 & Kontrolní součet & Výpočet SHA-256 pro pozdější ověření neporušenosti \\
7 & Údržba retence & Odstranění bodů obnovy nad rámec retenčního plánu \\
8 & Signalizace & Potvrzení úspěchu monitoringu; při selhání e-mailové upozornění \\
\bottomrule
\end{tabularx}
\end{center}

\subsection{Co přesně se zálohuje}

\subsubsection*{Databáze}

Výpis se pořizuje oficiálním nástrojem Supabase CLI, doporučeným postupem rozděleným do tří
souborů~\cite{ciautomat}. Toto rozdělení umožňuje obnovit strukturu bez dat, data bez struktury
nebo obojí --- podle povahy havárie.

\begin{center}
\begin{tabularx}{\textwidth}{@{}p{2.6cm}p{4.6cm}X@{}}
\toprule
\textbf{Soubor} & \textbf{Obsah} & \textbf{K čemu slouží} \\
\midrule
\texttt{roles.sql} & Databázové role a~oprávnění & Obnova přístupových práv \\
\texttt{schema.sql} & Struktura tabulek, indexy, pravidla RLS, pohledy, funkce & Postavení prázdné, ale funkční kopie \\
\texttt{data.sql} & Vlastní obsah tabulek & Naplnění daty \\
\bottomrule
\end{tabularx}
\end{center}

\subsubsection*{Nahrané soubory}

Úložiště souborů zpřístupňuje rozhraní kompatibilní s~S3~\cite{s3kompatibilita}, což umožňuje
použít standardní nástroj \texttt{rclone}. Synchronizace je \textbf{přírůstková} --- porovnává se
kontrolní součet a~přenášejí se pouze nové a~změněné soubory. Druhý a~každý další běh je proto
řádově rychlejší a~nezatěžuje datový přenos.

\begin{upozorneni}[Toto je jediná ochrana nahraných souborů, která existuje]
Vzhledem k~omezení popsanému v~kapitole~\ref{sec:storage-nezalohy} je tato část vrstvy 2
\textbf{jedinou} zálohou fotografií a~dokumentů v~celém systému. Pokud by se škola rozhodla
vrstvu 2 nezavést, zůstanou nahrané soubory zcela bez zálohy.
\end{upozorneni}

\subsection{Frekvence a~retenční plán}

Zálohovací cyklus běží \textbf{jednou týdně} pro databázi i~soubory zároveň. Retence sleduje
model \emph{děd -- otec -- syn}, který drží hustou historii pro nedávné události a~řídkou pro
vzdálené:

\begin{center}
\begin{tabularx}{\textwidth}{@{}p{3.4cm}p{2.8cm}Xc@{}}
\toprule
\textbf{Úroveň} & \textbf{Perioda} & \textbf{Pokrývá} & \textbf{Počet} \\
\midrule
Týdenní & každý týden & poslední dva měsíce & 8 \\
Měsíční & poslední v~měsíci & poslední rok & 12 \\
Roční & poslední v~roce & dlouhodobá archivace & 3 \\
\midrule
\multicolumn{3}{@{}l}{\textbf{Celkem bodů obnovy}} & \textbf{23} \\
\bottomrule
\end{tabularx}
\end{center}

\medskip
Tento plán znamená, že problém lze odhalit a~napravit i~po měsících. Roční body obnovy zároveň
umožňují doložit stav dat na konci každého akademického roku.

\begin{informace}[Poznámka k~RPO vrstvy 2]
Týdenní frekvence znamená, že vrstva 2 může ztratit až sedm dní práce. To je vědomý kompromis:
vrstva 2 \textbf{není} první linií obrany --- tou je vrstva 1 s~denní frekvencí. Vrstva 2 řeší
scénáře, které vrstva 1 nepokrývá vůbec (soubory, delší horizont, nezávislost na dodavateli),
a~pro ty je týdenní perioda dostatečná. Pokud by se v~budoucnu ukázala jako nedostatečná, lze
frekvenci zvýšit změnou jediného řádku konfigurace, bez zásahu do architektury.
\end{informace}

\subsection{Kalkulace potřebné kapacity}
\label{sec:kapacita}

Tarif Pro zahrnuje 100~GB prostoru pro soubory a~8~GB diskového prostoru databáze
(s~možností navýšení)~\cite{cenik}. Kalkulace níže počítá s~\textbf{plným vyčerpáním} zahrnutých
kapacit, tedy s~cílovým, nikoli současným stavem --- aby alokace nebyla za rok znovu předmětem
jednání.

\begin{center}
\begin{tabularx}{\textwidth}{@{}Xrrr@{}}
\toprule
\textbf{Položka} & \textbf{Jednotka} & \textbf{Počet} & \textbf{Celkem} \\
\midrule
Aktuální zrcadlo souborů & 100 GB & 1 & 100 GB \\
Verzované změny souborů (5\,\%/měsíc, 12 měsíců) & 5 GB & 12 & 60 GB \\
Týdenní výpisy databáze (komprimované) & 5 GB & 8 & 40 GB \\
Měsíční výpisy databáze & 5 GB & 12 & 60 GB \\
Roční výpisy databáze & 5 GB & 3 & 15 GB \\
Pracovní prostor pro rozbalení a~ověření obnovy & --- & --- & 150 GB \\
\midrule
\textbf{Součet} & & & \textbf{425 GB} \\
Rezerva na růst dat a~prodloužení retence (přibližně 2$\times$) & & & 425 GB \\
\midrule
\textbf{Navrhovaná alokace} & & & \textbf{1 TB} \\
\bottomrule
\end{tabularx}
\end{center}

\medskip
Rezerva není nadsazená: zaručuje, že zálohovací systém nepřestane fungovat kvůli zaplnění disku
--- což je nejčastější tichá příčina selhání zálohování. Pracovní prostor je nutný proto, že
ověření obnovy vyžaduje archiv skutečně rozbalit a~načíst do dočasné databáze.

\subsection{Zabezpečení záloh}
\label{sec:zabezpeceni}

Záloha obsahuje kompletní osobní údaje a~autorské práce studujících. Její zabezpečení musí být
proto přinejmenším na úrovni produkčního systému.

\begin{center}
\begin{tabularx}{\textwidth}{@{}p{4.4cm}X@{}}
\toprule
\textbf{Opatření} & \textbf{Provedení} \\
\midrule
Šifrování v~klidu & Každý archiv je zašifrován asymetricky. Server zná pouze \emph{veřejný} klíč, kterým šifruje. Bez soukromého klíče je obsah nečitelný i~pro správu školního serveru. \\
Úschova klíče & Soukromý klíč drží \textbf{dvě nezávislé strany}: technické vedení Tappky a~pověřená osoba školy. Ani jedna strana nemůže data zpřístupnit ani znepřístupnit sama. \\
Kontrolní součty & Ke každému archivu je uložen otisk SHA-256, který umožňuje kdykoli ověřit, že soubor nebyl poškozen ani pozměněn. \\
Neměnitelnost & Adresář s~archivy je nastaven jen pro zápis; zálohovací proces nemá právo mazat, mazání provádí oddělený proces správy retence. \\
Přístupová práva & Archivy jsou čitelné pouze pro určený systémový účet. \\
Tajemství & Přístupové údaje k~Supabase jsou uloženy v~souboru mimo obraz kontejneru, s~právy pro čtení jen vlastníkem. \\
\bottomrule
\end{tabularx}
\end{center}

\begin{informace}[Proč dvojí úschova klíče]
Rozdělení klíče mezi školu a~vývojový tým řeší dvě protichůdná rizika najednou. Správa školního
serveru \textbf{nemá} důvod číst eseje a~reflexe studujících --- šifrování jí v~tom brání. Zároveň
škola \textbf{nesmí} být závislá na existenci vývojového týmu --- proto drží vlastní kopii klíče
a~data si dokáže zpřístupnit i~v~případě, že by spolupráce s~týmem skončila.
\end{informace}

\subsection{Monitoring a~ověřování obnovitelnosti}
\label{sec:overovani}

Nejběžnější způsob, jak zálohovací systém selže, není havárie --- je to \emph{tiché} přerušení
běhu, které si nikdo měsíce nevšimne. Návrh proto obsahuje aktivní kontrolu:

\begin{center}
\begin{tabularx}{\textwidth}{@{}p{3.6cm}p{2.6cm}X@{}}
\toprule
\textbf{Kontrola} & \textbf{Perioda} & \textbf{Popis} \\
\midrule
Signál o~dokončení & každý běh & Po úspěšném cyklu je odeslán potvrzovací signál. \textbf{Chybějící} signál vyvolá upozornění --- hlídá se tedy i~případ, kdy se úloha vůbec nespustí. \\
Kontrola velikosti & každý běh & Výpisy nesmí být prázdné ani výrazně menší než minule. \\
Ověření otisku & měsíčně & Přepočet SHA-256 náhodně zvoleného archivu. \\
\textbf{Zkušební obnova} & čtvrtletně & Archiv se rozbalí a~načte do dočasné databáze; ověří se počty záznamů a~funkčnost. Výsledek se zaznamená do protokolu. \\
Úplné nácvikové obnovení & ročně & Postavení plně funkční kopie Tappky ze zálohy, včetně souborů. \\
\bottomrule
\end{tabularx}
\end{center}

\medskip
Čtvrtletní zkušební obnova je to, co odlišuje skutečnou zálohu od souboru, o~kterém se
\emph{předpokládá}, že je zálohou. Odborná doporučení uvádějí, že samotná existence archivu není
důkazem obnovitelnosti --- ověřením je až úspěšné načtení dat a~kontrola jejich
konzistence~\cite{pravidlo321}.

\subsection{Proces vyžádání dat}
\label{sec:vyzadani}

Záloha má hodnotu jen tehdy, je-li dostupná ve chvíli, kdy je potřeba. Navrhujeme proto předem
dohodnout následující lhůty. Bez nich hrozí, že v~krizové situaci bude nutné teprve hledat, kdo
má oprávnění a~koho kontaktovat.

\begin{center}
\begin{tabularx}{\textwidth}{@{}p{4.2cm}p{3.4cm}X@{}}
\toprule
\textbf{Typ žádosti} & \textbf{Lhůta} & \textbf{Popis} \\
\midrule
\textbf{Havarijní obnova} & \textbf{do 24 hodin} & Výpadek nebo ztráta dat na straně Supabase. Předání nejnovějšího bodu obnovy má absolutní prioritu; mimo pracovní dobu dle předem sjednaného kontaktu. \\
Předání jednoho bodu obnovy & do 2 pracovních dnů & Běžná situace --- obnova omylem smazaných dat či souborů. \\
Předání celé sady archivů & do 5 pracovních dnů & Migrace, audit, ukončení nebo změna provozu. \\
Ověření existence záloh & do 2 pracovních dnů & Doložení, že zálohovací proces běží (protokol posledních běhů). \\
\bottomrule
\end{tabularx}
\end{center}

\subsubsection*{Postup}

\begin{enumerate}[leftmargin=1.6em]
  \item Vývojový tým Tappky zašle písemnou žádost pověřenému kontaktu oddělení PEVKA s~uvedením
        požadovaného bodu obnovy a~důvodu.
  \item Oddělení PEVKA ověří oprávněnost žádosti a~zpřístupní zašifrované archivy dohodnutým
        zabezpečeným kanálem (přenos v~rámci školní sítě nebo předání fyzického média).
  \item Příjemce ověří kontrolní součet SHA-256 a~teprve poté archiv dešifruje soukromým klíčem.
  \item O~předání se pořídí stručný záznam --- kdo, kdy, co a~z~jakého důvodu. Záznam slouží jako
        doklad o~zacházení s~osobními údaji.
\end{enumerate}

\begin{informace}[Alternativa k~procesu vyžádání]
Pokud by lhůty byly z~provozních důvodů problematické, lze zvolit variantu, kdy má vývojový tým
\textbf{přímý přístup pro čtení} k~adresáři se zálohami (například přes SFTP s~omezenými právy).
Archivy jsou šifrované, takže tento přístup sám o~sobě nezvyšuje riziko úniku dat --- pouze
zkracuje reakční dobu na minuty. Volba mezi oběma variantami je na uvážení školy.
\end{informace}

\subsection{Souhrn požadavků na školu}

\begin{center}
\begin{tabularx}{\textwidth}{@{}p{0.7cm}p{3.6cm}X@{}}
\toprule
& \textbf{Oblast} & \textbf{Požadavek} \\
\midrule
$\square$ & Server & Virtuální nebo fyzický stroj, Linux (doporučeno Debian 12 nebo Ubuntu 24.04 LTS) \\
$\square$ & Výkon & 2 vCPU, 4 GB RAM --- úloha je I/O náročná, nikoli výpočetně \\
$\square$ & Úložiště & \textbf{1 TB} dedikovaného prostoru; ideálně na jiném fyzickém poli než produkční systémy školy \\
$\square$ & Software & Docker Engine (verze 24 a~vyšší) s~pluginem Docker Compose \\
$\square$ & Síť --- odchozí & TCP 443 (HTTPS) na \texttt{*.supabase.co} a~\texttt{*.supabase.com} \\
$\square$ & Síť --- odchozí & TCP 5432 na databázový server Supabase (výpis databáze) \\
$\square$ & Síť --- příchozí & Žádný požadavek; služba nevystavuje žádný port \\
$\square$ & Notifikace & E-mailový relay nebo webhook pro upozornění na selhání \\
$\square$ & Přístup & Účet pro prvotní instalaci a~čtvrtletní ověření obnovy \\
$\square$ & Kontakt & Pověřená osoba oddělení PEVKA pro proces dle kap.~\ref{sec:vyzadani} \\
$\square$ & Úschova klíče & Určení osoby, která bude držet školní kopii soukromého klíče \\
\bottomrule
\end{tabularx}
\end{center}

\medskip
Zátěž serveru je nízká a~soustředěná do jednoho okna v~týdnu. Mimo běh zálohy je spotřeba
prostředků prakticky nulová --- stroj lze proto sdílet s~jinými nenáročnými službami.

\newpage

% =========================================================================
\section{Vrstva 3 --- lidsky čitelné exporty (koncept)}
\label{sec:vrstva3}

Vrstvy 1 a~2 chrání data v~\emph{technické} podobě. Obě předpokládají, že v~okamžiku potřeby
bude k~dispozici někdo, kdo umí obnovit databázi a~spustit aplikaci. Vrstva 3 řeší otázku, která
tímto předpokladem není pokryta:

\begin{center}
\emph{Co když je potřeba se v~datech vyznat bez Tappky a~bez technických znalostí?}
\end{center}

\subsection{Myšlenka}

Umožnit přímo v~rozhraní Tappky vyexportovat data do \textbf{běžných, čitelných formátů} ---
typicky jako archiv ZIP obsahující dokumenty, tabulky a~nahrané soubory, uspořádané do
srozumitelné adresářové struktury. Takový archiv by šlo otevřít na jakémkoli počítači, prohlédnout
v~běžném prohlížeči souborů a~přečíst bez jakéhokoli speciálního nástroje.

Export by dával smysl pořizovat \textbf{na konci každého semestru} a~archivovat spolu se zálohami
vrstvy 2. Zahrnoval by texty i~přiložené soubory --- na rozdíl od databázového výpisu, který je
sice úplný, ale bez aplikace prakticky nečitelný.

\subsection{Jaké situace by řešil}

\begin{itemize}
  \item Provoz Tappky by z~jakéhokoli důvodu skončil a~škola potřebuje záznamy zachovat
        v~použitelné podobě.
  \item Vyučující potřebuje nahlédnout do materiálů v~době výpadku aplikace.
  \item Studující žádá o~kopii vlastních dat (právo na přenositelnost údajů dle GDPR).
  \item Škola potřebuje doložit stav dokumentace za uzavřený akademický rok nezávisle na
        běžícím systému.
\end{itemize}

\begin{informace}[Stav]
Vrstva 3 je v~tuto chvíli \textbf{koncept}, nikoli hotová funkce. Není součástí požadavků na
školu a~nevyžaduje žádné rozhodnutí --- je zde uvedena proto, aby byl model ochrany dat popsán
celý a~aby bylo zřejmé, jakým směrem se má do budoucna rozvíjet. Konkrétní podoba exportu bude
předmětem samostatného návrhu.
\end{informace}

\newpage

% =========================================================================
\section{Souhrnné srovnání vrstev}

\begin{center}
\begin{tabularx}{\textwidth}{@{}p{2.7cm}XXX@{}}
\toprule
 & \textbf{Vrstva 1} & \textbf{Vrstva 2} & \textbf{Vrstva 3} \\
 & {\footnotesize Supabase Pro} & {\footnotesize školní server} & {\footnotesize čitelné exporty} \\
\midrule
Vlastník & Supabase & \textbf{Škola} & Škola \\
Umístění & Zurich, \texttt{eu-central-2} & Česká republika & Dle volby školy \\
Frekvence & denně & týdně & semestrálně \\
Retence & 7 dní & 23 bodů obnovy (až 3 roky) & neomezená \\
Databáze & \ano & \ano & částečně \\
\textbf{Nahrané soubory} & \textbf{\nepokr} & \textbf{\ano} & \ano \\
RPO & do 24 h & do 7 dní & do 1 semestru \\
RTO & desítky minut až hodiny & 1--2 pracovní dny & okamžitě \\
Čitelné bez aplikace & \nepokr & \nepokr & \ano \\
Nezávislé na dodavateli & \nepokr & \ano & \ano \\
Náklady & součást tarifu & serverová kapacita & vývojová kapacita \\
Stav & \ano v~provozu & \textbf{předmět tohoto návrhu} & koncept \\
\bottomrule
\end{tabularx}
\end{center}

\medskip
\noindent
\textbf{RPO} (\emph{Recovery Point Objective}) --- maximální množství práce, o~které lze v~nejhorším
případě přijít, vyjádřené v~čase.\\
\textbf{RTO} (\emph{Recovery Time Objective}) --- doba od rozhodnutí o~obnově do okamžiku, kdy jsou
data opět k~dispozici.

% =========================================================================
\section{Doporučený postup}

\begin{center}
\begin{tabularx}{\textwidth}{@{}p{0.7cm}p{4.6cm}Xp{2.4cm}@{}}
\toprule
\textbf{\#} & \textbf{Krok} & \textbf{Popis} & \textbf{Odpovídá} \\
\midrule
1 & Projednání a~schválení & Odsouhlasení alokace serveru a~1 TB prostoru & Vedení školy \\
2 & Přidělení serveru & Zřízení stroje, instalace Dockeru, otevření odchozích portů & IT oddělení \\
3 & Generování klíčů & Vytvoření šifrovacího páru, předání soukromého klíče oběma stranám & Tappka $+$ škola \\
4 & Nasazení & Nahrání konfigurace, první ruční spuštění, ověření výsledku & Tappka $+$ IT \\
5 & Ověření obnovy & První zkušební obnova a~sepsání protokolu & Tappka \\
6 & Předání & Předání provozní dokumentace a~kontaktů & Tappka \\
7 & Dohoda o~lhůtách & Písemné potvrzení lhůt dle kap.~\ref{sec:vyzadani} & Škola $+$ PEVKA \\
8 & Rutinní provoz & Čtvrtletní ověření obnovy, roční nácvik & Tappka \\
\bottomrule
\end{tabularx}
\end{center}

\medskip
Kroky 2 až 6 představují řádově jeden až dva dny práce. Následný provoz je automatický.

\newpage

% =========================================================================
\appendix
\section{Příloha: Konfigurace služby}

Následující výpisy jsou ilustrativní --- ukazují rozsah a~povahu konfigurace, aby si IT oddělení
mohlo udělat představu o~nárocích. Finální podoba vznikne při nasazení.

\subsection{Definice služby (\texttt{docker-compose.yml})}

\begin{lstlisting}[language=bash]
services:
  tappka-backup:
    image: ghcr.io/tappka/backup:1
    container_name: tappka-backup
    restart: unless-stopped
    env_file: /etc/tappka/backup.env
    environment:
      # Every Sunday at 02:00 UTC
      BACKUP_CRON: "0 2 * * 0"
      BUCKETS: "avatars images documents"
    volumes:
      - /srv/tappka-backup:/backup
    logging:
      driver: json-file
      options: { max-size: "10m", max-file: "5" }
\end{lstlisting}

\subsection{Tajemství (\texttt{/etc/tappka/backup.env}, práva 0600)}

\begin{lstlisting}[language=bash]
# PostgreSQL connection string (Supabase project)
SUPABASE_DB_URL=postgresql://...

# S3-compatible endpoint for Storage buckets
RCLONE_CONFIG_SUPABASE_TYPE=s3
RCLONE_CONFIG_SUPABASE_PROVIDER=Other
RCLONE_CONFIG_SUPABASE_ENDPOINT=https://<project>.supabase.co/storage/v1/s3
RCLONE_CONFIG_SUPABASE_ACCESS_KEY_ID=...
RCLONE_CONFIG_SUPABASE_SECRET_ACCESS_KEY=...

# Public key used to encrypt archives (private key is NOT stored here)
AGE_RECIPIENT=age1...

# Monitoring: a missing ping raises an alert
HEALTHCHECK_URL=https://...
\end{lstlisting}

\subsection{Zálohovací skript (zkráceně)}

\begin{lstlisting}[language=bash]
#!/bin/sh
set -eu

STAMP=$(date -u +%Y-%m-%dT%H%M%SZ)
WORK="/backup/work/${STAMP}"
OUT="/backup/archive"
mkdir -p "$WORK" "$OUT"

# --- 1. Database: roles, schema, data (three separate files) -------------
supabase db dump --db-url "$SUPABASE_DB_URL" -f "$WORK/roles.sql"  --role-only
supabase db dump --db-url "$SUPABASE_DB_URL" -f "$WORK/schema.sql"
supabase db dump --db-url "$SUPABASE_DB_URL" -f "$WORK/data.sql" \
    --data-only --use-copy

# --- 2. Sanity check: refuse to keep an empty or truncated dump ----------
for f in roles schema data; do
  test -s "$WORK/$f.sql" || { echo "FATAL: $f.sql is empty"; exit 1; }
done

# --- 3. Storage buckets: incremental sync, changes kept as versions ------
for BUCKET in $BUCKETS; do
  rclone sync "supabase:${BUCKET}" "/backup/storage/${BUCKET}" \
    --backup-dir "/backup/storage-versions/${STAMP}/${BUCKET}" \
    --checksum --transfers 8 --retries 5
done

# --- 4. Pack, encrypt, checksum -----------------------------------------
tar -czf - -C "$WORK" . \
  | age -r "$AGE_RECIPIENT" -o "$OUT/tappka-db-${STAMP}.tar.gz.age"
sha256sum "$OUT/tappka-db-${STAMP}.tar.gz.age" \
  > "$OUT/tappka-db-${STAMP}.sha256"
rm -rf "$WORK"

# --- 5. Retention + heartbeat -------------------------------------------
/usr/local/bin/prune-gfs.sh "$OUT"
curl -fsS "$HEALTHCHECK_URL"
\end{lstlisting}

\section{Příloha: Manuál spuštění}

\begin{enumerate}[leftmargin=1.6em]
  \item \textbf{Příprava serveru.} Instalace Docker Engine s~pluginem Compose. Vytvoření adresáře
        \texttt{/srv/tappka-backup} na svazku s~1 TB volného prostoru.
  \item \textbf{Šifrovací klíč.} Vygenerování páru klíčů. Veřejný klíč se vloží do konfigurace na
        serveru; soukromý klíč se \emph{na server neukládá} a~předá se zabezpečeným kanálem
        oběma stranám, které jej drží (viz kap.~\ref{sec:zabezpeceni}).
  \item \textbf{Přístupové údaje.} Vytvoření souboru \texttt{/etc/tappka/backup.env} s~právy 0600
        a~vlastníkem \texttt{root}.
  \item \textbf{Ověření sítě.} Test dostupnosti databázového serveru a~koncového bodu úložiště
        ze serveru školy.
  \item \textbf{První ruční běh.}
        \begin{lstlisting}[language=bash]
docker compose run --rm tappka-backup /usr/local/bin/backup.sh
        \end{lstlisting}
  \item \textbf{Kontrola výsledku.} Ověření, že v~\texttt{/backup/archive} vznikl archiv nenulové
        velikosti a~odpovídající soubor s~kontrolním součtem.
  \item \textbf{Trvalé spuštění.}
        \begin{lstlisting}[language=bash]
docker compose up -d
        \end{lstlisting}
  \item \textbf{Ověření monitoringu.} Záměrné vynechání jednoho běhu a~kontrola, že dorazí
        upozornění.
  \item \textbf{Zkušební obnova.} Provedení postupu z~přílohy~\ref{sec:obnova} a~sepsání protokolu.
\end{enumerate}

\section{Příloha: Postup obnovy}
\label{sec:obnova}

\begin{enumerate}[leftmargin=1.6em]
  \item \textbf{Volba bodu obnovy.} Určení, ke kterému datu se obnovuje.
  \item \textbf{Ověření neporušenosti.} Přepočet SHA-256 a~porovnání s~uloženým otiskem.
        \emph{Neshoda znamená, že archiv nelze použít --- pokračuje se předchozím bodem obnovy.}
  \item \textbf{Dešifrování} soukromým klíčem a~rozbalení archivu.
  \item \textbf{Obnova databáze} v~pořadí \texttt{roles.sql} $\rightarrow$ \texttt{schema.sql}
        $\rightarrow$ \texttt{data.sql}.
  \item \textbf{Obnova souborů} nahráním zrcadla kontejnerů \texttt{avatars}, \texttt{images}
        a~\texttt{documents} zpět do úložiště.
  \item \textbf{Kontrola konzistence.} Porovnání počtu záznamů, namátková kontrola, že odkazy na
        soubory vedou na existující objekty.
  \item \textbf{Ruční kroky po obnově.} Obnovení hesel vlastních databázových rolí a~případných
        replikačních odběrů --- ty nejsou součástí výpisu.
  \item \textbf{Funkční ověření.} Přihlášení do aplikace, otevření eseje s~přílohou, zobrazení
        profilové fotografie.
  \item \textbf{Protokol.} Zaznamenání data, použitého bodu obnovy, doby trvání a~zjištěných
        odchylek.
\end{enumerate}

\newpage

% =========================================================================
\section{Zdroje}
\label{sec:zdroje}

Všechny odkazy ověřeny k~31.\,8.\,2026.

\begin{thebibliography}{9}
\small

\bibitem{zalohysupabase}
\textsc{Supabase.} \emph{Database Backups}. Oficiální dokumentace.\\
\url{https://supabase.com/docs/guides/platform/backups}\\
{\color{muted}Doloženo: retence 7 dní denních záloh v~tarifu Pro; fyzické zálohy; zálohy
nezahrnují objekty uložené přes Storage API; nedostupnost projektu po dobu obnovy.}

\bibitem{cenik}
\textsc{Supabase.} \emph{Pricing}.\\
\url{https://supabase.com/pricing}\\
{\color{muted}Doloženo: obsah tarifu Pro --- 8 GB diskového prostoru databáze, 100 GB prostoru
pro soubory, 7denní retence logů, denní zálohy po 7 dní, absence SLA, SOC 2, ISO 27001 a~SSO.}

\bibitem{sdilenaodpovednost}
\textsc{Supabase.} \emph{Shared Responsibility Model}.\\
\url{https://supabase.com/docs/guides/deployment/shared-responsibility-model}\\
{\color{muted}Doloženo: rozdělení odpovědnosti mezi Supabase a~provozovatele aplikace; citovaná
věta o~minimu zábran kolem databáze.}

\bibitem{soc2}
\textsc{Supabase.} \emph{SOC 2 Compliance and Supabase}.\\
\url{https://supabase.com/docs/guides/security/soc-2-compliance}\\
{\color{muted}Doloženo: Supabase je SOC 2 Type 2 compliant a~je každoročně auditována; auditní
zpráva je dostupná pouze v~tarifech Team a~Enterprise.}

\bibitem{prodchecklist}
\textsc{Supabase.} \emph{Production Checklist}.\\
\url{https://supabase.com/docs/guides/platform/going-into-prod}\\
{\color{muted}Doloženo: dostupnost vynucení TLS, síťových omezení, RLS, vícefaktorového ověření
a~omezení četnosti požadavků.}

\bibitem{ciautomat}
\textsc{Supabase.} \emph{Automated backups using GitHub Actions}.\\
\url{https://supabase.com/docs/guides/deployment/ci/backups}\\
{\color{muted}Doloženo: doporučený postup výpisu do tří souborů (\texttt{roles.sql},
\texttt{schema.sql}, \texttt{data.sql}) včetně přepínačů \texttt{--role-only},
\texttt{--data-only} a~\texttt{--use-copy}.}

\bibitem{s3kompatibilita}
\textsc{Supabase.} \emph{S3 Compatibility}.\\
\url{https://supabase.com/docs/guides/storage/s3/compatibility}\\
{\color{muted}Doloženo: úložiště souborů zpřístupňuje rozhraní kompatibilní s~S3 na koncovém bodu
\texttt{/storage/v1/s3}, použitelné standardními nástroji.}

\bibitem{adekvatnost}
\textsc{Federální úřad pro ochranu údajů a~transparentnost (EDÖB, Švýcarsko).}
\emph{EU adequacy decision regarding Switzerland}, 15.\,1.\,2024.\\
\url{https://www.edoeb.admin.ch/en/15012024-eu-adequacy-decision-regarding-switzerland}\\
{\color{muted}Doloženo: Evropská komise potvrdila dne 15.\,1.\,2024 odpovídající úroveň ochrany
osobních údajů ve Švýcarsku; přenos z~EU/EHP je možný bez dodatečných záruk.}

\bibitem{pravidlo321}
\textsc{Přehled odborných doporučení k~pravidlu 3--2--1} (Veeam, Commvault, Object First).\\
\url{https://www.veeam.com/blog/321-backup-rule.html}\\
{\color{muted}Doloženo: pravidlo 3--2--1 jako základní standard doporučovaný CISA a~NIST;
rozšíření 3--2--1--1--0 o~neměnitelnou kopii a~ověřované obnovy; nutnost testovat obnovu na
úrovni aplikace, nikoli pouze existenci archivu.}

\end{thebibliography}

\end{document}
